\documentclass[aspectratio=169,10pt]{beamer}
% ---------------------------------------------------------------
% R4.02 Qualite de developpement -- preambule commun des diapositives
% Charte reprise de R3.07 (metropolis, dark tech)
% Palette IUT partagee (violet du logo + accent or)
% Utilisation :
%   \documentclass[aspectratio=169,10pt]{beamer}
%   % ---------------------------------------------------------------
% R4.02 Qualite de developpement -- preambule commun des diapositives
% Charte reprise de R3.07 (metropolis, dark tech)
% Palette IUT partagee (violet du logo + accent or)
% Utilisation :
%   \documentclass[aspectratio=169,10pt]{beamer}
%   % ---------------------------------------------------------------
% R4.02 Qualite de developpement -- preambule commun des diapositives
% Charte reprise de R3.07 (metropolis, dark tech)
% Palette IUT partagee (violet du logo + accent or)
% Utilisation :
%   \documentclass[aspectratio=169,10pt]{beamer}
%   \input{../../latex/common/r402-slides-preamble.tex}
%   \graphicspath{{../../latex/common/assets/}{img/}{figs/}}
% ---------------------------------------------------------------
\usetheme{metropolis}
\usepackage{appendixnumberbeamer}

\usepackage[T1]{fontenc}
\usepackage[utf8]{inputenc}
\usepackage[french]{babel}
\usepackage[scale=0.94]{tgheros}
\renewcommand{\familydefault}{\sfdefault}
\usepackage{amsmath,amssymb,amsfonts}
\usepackage{graphicx}
\usepackage{booktabs}
\usepackage{multirow}
\usepackage{tikz}
\usetikzlibrary{shapes.geometric, shapes, arrows, arrows.meta, positioning, calc, fit}
\usepackage{xcolor}
\usepackage{listings}
\usepackage[skins,listings]{tcolorbox}
\usepackage{qrcode}
\usepackage{eurosym}
\usepackage{hyperref}

% ---------------------------------------------------------------
% PALETTE (dark-mode tech: near-black canvas, violet (logo IUT) + gold accents)
% ---------------------------------------------------------------
\definecolor{qabg}{HTML}{0E1016}      % canvas background
\definecolor{qapanel}{HTML}{181B24}   % card / block panel
\definecolor{qapanelalt}{HTML}{20242F}% secondary panel tint
\definecolor{qaline}{HTML}{2B303D}    % hairline rules / table rules
\definecolor{qatext}{HTML}{ECEEF2}    % primary text (off-white)
\definecolor{qamuted}{HTML}{8B93A6}   % secondary / caption text
\definecolor{qaaccent}{HTML}{D93ED7}  % primary accent -- violet du logo IUT, eclairci pour fond sombre
\definecolor{qaaccent2}{HTML}{F2B84B} % secondary accent -- or/ambre
\definecolor{qacodekw}{HTML}{C792EA}  % code : mots-cles
\definecolor{qacodestr}{HTML}{C3E88D} % code : chaines
\definecolor{qacodecom}{HTML}{7F8AA3} % code : commentaires
\definecolor{qacodeann}{HTML}{F2B84B} % code : annotations / javadoc

\setbeamercolor{background canvas}{bg=qabg}
\setbeamercolor{normal text}{fg=qatext}
\setbeamercolor{palette primary}{bg=qabg, fg=qamuted}
\setbeamercolor{frametitle}{bg=qabg, fg=qatext}
\setbeamercolor{progress bar}{fg=qaaccent, bg=qaline}
\setbeamercolor{alerted text}{fg=qaaccent2}
\setbeamercolor{block title}{bg=qapanel, fg=qaaccent}
\setbeamercolor{block body}{bg=qapanel, fg=qatext}
\setbeamercolor{block title alerted}{bg=qapanel, fg=qaaccent2}
\setbeamercolor{block body alerted}{bg=qapanel, fg=qatext}
\setbeamercolor{block title example}{bg=qapanelalt, fg=qaaccent}
\setbeamercolor{block body example}{bg=qapanelalt, fg=qatext}
\setbeamercolor{example text}{fg=qaaccent}
\setbeamercolor{itemize item}{fg=qaaccent}
\setbeamercolor{itemize subitem}{fg=qaaccent2}
\setbeamercolor{enumerate item}{fg=qaaccent}
\setbeamercolor{enumerate subitem}{fg=qaaccent2}
\setbeamercolor{caption}{fg=qamuted}
\setbeamercolor{page number in head/foot}{fg=qamuted, bg=qabg}
\setbeamercolor{title}{fg=qatext}
\setbeamercolor{subtitle}{fg=qamuted}
\setbeamercolor{author}{fg=qatext}
\setbeamercolor{institute}{fg=qamuted}
\setbeamercolor{date}{fg=qamuted}
\setbeamercolor{section in toc}{fg=qaaccent}
\setbeamercolor{subsection in toc}{fg=qamuted}
% Sommaire de debut de section : section en cours en violet et gras, les autres en gris
\setbeamercolor{section in toc shaded}{fg=qamuted}
\setbeamercolor{subsection in toc shaded}{fg=qamuted}

\setbeamertemplate{blocks}[rounded][shadow=false]
\setbeamertemplate{itemize item}{\raisebox{0.15em}{\scriptsize\textcolor{qaaccent}{\ensuremath{\blacktriangleright}}}}
\setbeamertemplate{itemize subitem}{\raisebox{0.1em}{\tiny\textcolor{qaaccent2}{\ensuremath{\bullet}}}}

\setbeamerfont{frametitle}{series=\bfseries, size=\Large}
\setbeamerfont{title}{series=\bfseries, size=\Large}
\setbeamerfont{subtitle}{size=\normalsize}
\setbeamerfont{institute}{size=\scriptsize}
\setbeamerfont{date}{size=\scriptsize}
\setbeamerfont{author}{size=\small}

\metroset{sectionpage=none, subsectionpage=none, numbering=fraction, progressbar=frametitle}
\setbeamertemplate{section in toc}[sections numbered]

% ---------------------------------------------------------------
% LISTINGS -- theme sombre (lstlisting reste utilisable tel quel)
% \lstcode : police + couleur de base, a reutiliser dans les
%            surcharges locales (ex. basicstyle=\lstcode\scriptsize)
% ---------------------------------------------------------------
\newcommand{\lstcode}{\ttfamily\color{qatext}}

\lstdefinelanguage{TypeScript}{
  keywords={typeof, new, true, false, catch, function, return, null, switch, var,
    if, in, while, do, else, case, break, const, let, interface, class, extends,
    implements, enum, export, import, this, super, constructor, public, private,
    protected, static, async, await, throw, throws, try, finally, void},
  ndkeywords={boolean, any, number, string, int, long, double},
  sensitive=true,
  comment=[l]{//},
  morecomment=[s]{/*}{*/},
  morestring=[b]',
  morestring=[b]"
}

\lstset{
  language=Java,
  basicstyle=\lstcode\footnotesize,
  keywordstyle=\color{qacodekw}\bfseries,
  ndkeywordstyle=\color{qaaccent2}\bfseries,
  identifierstyle=\color{qatext},
  stringstyle=\color{qacodestr},
  commentstyle=\color{qacodecom}\itshape,
  morecomment=[s][\color{qacodeann}]{/**}{*/},
  backgroundcolor=\color{qapanelalt},
  rulecolor=\color{qaline},
  frame=single,
  framesep=3pt,
  xleftmargin=4pt,
  xrightmargin=4pt,
  keepspaces=true,
  columns=fullflexible,
  showstringspaces=false,
  breaklines=true,
  tabsize=2,
  escapeinside={(*@}{@*)},
  % accents UTF-8 dans le code (listings ne gere pas l'UTF-8 multi-octets)
  literate={é}{{\'e}}1 {è}{{\`e}}1 {ê}{{\^e}}1 {à}{{\`a}}1 {â}{{\^a}}1
           {ù}{{\`u}}1 {û}{{\^u}}1 {ô}{{\^o}}1 {î}{{\^i}}1 {ç}{{\c c}}1
           {É}{{\'E}}1 {À}{{\`A}}1 {€}{{\euro{}}}1 {…}{{\ldots}}1,
}

% ---------------------------------------------------------------
% CARD-STYLE BOXES
% ---------------------------------------------------------------
\tcbset{
  figcard/.style={colback=white, colframe=white, boxrule=0pt, arc=4pt,
    left=3pt, right=3pt, top=3pt, bottom=3pt, boxsep=0pt,
    halign=center, valign=center},
  datacard/.style={colback=qapanel, colframe=qaline, boxrule=0.4pt, arc=3pt,
    left=8pt, right=8pt, top=6pt, bottom=6pt, coltext=qatext,
    before upper={\color{qatext}}, halign=center},
  codebox/.style={colback=qapanelalt, colframe=qaline, boxrule=0.4pt, arc=3pt,
    left=8pt, right=8pt, top=5pt, bottom=5pt, listing only,
    listing options={basicstyle=\ttfamily\small\color{qaaccent},
      breaklines=true, showstringspaces=false, columns=fullflexible,
      keepspaces=true, frame=none, backgroundcolor={}}}
}
\newtcolorbox{figcard}[1][]{figcard,#1}
\newtcolorbox{datacard}{datacard}
\newtcblisting{codebox}{codebox}

% ---------------------------------------------------------------
% HANDY MACROS
% ---------------------------------------------------------------
% Mention d'elaboration (discrete : plus sombre que le texte secondaire)
\definecolor{qaaimention}{HTML}{4E5566}
\newcommand{\raimention}{{\tiny\color{qaaimention}Supports \'elabor\'es avec l'aide de Claude Code Opus 5.5}}
\newcommand{\hl}[1]{\alert{#1}}
\newcommand{\source}[1]{{\scriptsize\color{qamuted}#1}}

% Sommaire automatique en debut de section
\AtBeginSection[]{%
  \begin{frame}<beamer>{Plan du cours}
    \tableofcontents[currentsection]
  \end{frame}}

% Page de titre : \rtitleslide{numero du cours}{titre du cours}
% Les auteurs sont repris de \rauthors (modifiable avant \begin{document})
\newcommand{\rauthors}{Chouki Tibermacine}
\newcommand{\rtitleslide}[2]{%
\begin{frame}[plain,noframenumbering]
  \begin{center}
    \vspace{2mm}
    {\color{qaaccent}\rule{2.1cm}{1.2pt}}\\[6mm]
    {\Large\bfseries\color{qatext} R4.02}\\[1mm]
    {\Large\bfseries\color{qatext} Qualit\'e de d\'eveloppement}\\[5mm]
    {\large\color{qamuted} Cours #1}\\[1mm]
    {\large\color{qaaccent} #2}\\[6mm]
    \includegraphics[height=2.1cm]{iut_logo_transparent.png}

    \vspace{4mm}

    {\small\color{qatext} \rauthors}\\[1mm]
    {\scriptsize\color{qamuted} BUT 2 Informatique --- IUT de Vannes --- Universit\'e Bretagne Sud}\\[2mm]
    \raimention
  \end{center}
\end{frame}
}

% ---- Diapo de demonstration : \rdemo{numero}{titre}{contenu (items)}{fichier(s) a ouvrir} ----
\newcommand{\rdemo}[4]{%
\begin{frame}{D\'emo #1 --- #2}
  \begin{tikzpicture}[remember picture, overlay]
    \node[anchor=north east, fill=qaaccent2, text=qabg, font=\scriptsize\bfseries,
      rounded corners=3pt, inner sep=4pt] at ([xshift=-0.4cm, yshift=-0.35cm]current page.north east) {D\'EMO};
  \end{tikzpicture}
  \begin{columns}[T]
    \begin{column}{0.62\textwidth}
      \begin{itemize}
        #3
      \end{itemize}
    \end{column}
    \begin{column}{0.34\textwidth}
      \begin{datacard}
        \raggedright\scriptsize
        {\color{qamuted} Dans l'IDE}\\[1mm]
        #4
      \end{datacard}
    \end{column}
  \end{columns}
\end{frame}}

% ---- Diapo de fin « Questions ? » ----
\newcommand{\rcontact}{chouki.tibermacine@univ-ubs.fr}
\newcommand{\rquestionslide}{%
\begin{frame}{Questions ?}
  \begin{center}
    \vspace{2mm}
    \begin{tikzpicture}
      \node[regular polygon, regular polygon sides=6, draw=qaline, thick,
        minimum size=4.6cm, rotate=90] (hexin) {};
      \node[regular polygon, regular polygon sides=6, draw=qaaccent, dashed,
        thick, minimum size=6.2cm, rotate=90] (hexout) {};
      \node[font=\fontsize{92}{92}\selectfont\bfseries, color=qaaccent2] at (hexin) {?};
    \end{tikzpicture}\\[5mm]
    {\small\color{qamuted} \rcontact}
  \end{center}
\end{frame}}

%   \graphicspath{{../../latex/common/assets/}{img/}{figs/}}
% ---------------------------------------------------------------
\usetheme{metropolis}
\usepackage{appendixnumberbeamer}

\usepackage[T1]{fontenc}
\usepackage[utf8]{inputenc}
\usepackage[french]{babel}
\usepackage[scale=0.94]{tgheros}
\renewcommand{\familydefault}{\sfdefault}
\usepackage{amsmath,amssymb,amsfonts}
\usepackage{graphicx}
\usepackage{booktabs}
\usepackage{multirow}
\usepackage{tikz}
\usetikzlibrary{shapes.geometric, shapes, arrows, arrows.meta, positioning, calc, fit}
\usepackage{xcolor}
\usepackage{listings}
\usepackage[skins,listings]{tcolorbox}
\usepackage{qrcode}
\usepackage{eurosym}
\usepackage{hyperref}

% ---------------------------------------------------------------
% PALETTE (dark-mode tech: near-black canvas, violet (logo IUT) + gold accents)
% ---------------------------------------------------------------
\definecolor{qabg}{HTML}{0E1016}      % canvas background
\definecolor{qapanel}{HTML}{181B24}   % card / block panel
\definecolor{qapanelalt}{HTML}{20242F}% secondary panel tint
\definecolor{qaline}{HTML}{2B303D}    % hairline rules / table rules
\definecolor{qatext}{HTML}{ECEEF2}    % primary text (off-white)
\definecolor{qamuted}{HTML}{8B93A6}   % secondary / caption text
\definecolor{qaaccent}{HTML}{D93ED7}  % primary accent -- violet du logo IUT, eclairci pour fond sombre
\definecolor{qaaccent2}{HTML}{F2B84B} % secondary accent -- or/ambre
\definecolor{qacodekw}{HTML}{C792EA}  % code : mots-cles
\definecolor{qacodestr}{HTML}{C3E88D} % code : chaines
\definecolor{qacodecom}{HTML}{7F8AA3} % code : commentaires
\definecolor{qacodeann}{HTML}{F2B84B} % code : annotations / javadoc

\setbeamercolor{background canvas}{bg=qabg}
\setbeamercolor{normal text}{fg=qatext}
\setbeamercolor{palette primary}{bg=qabg, fg=qamuted}
\setbeamercolor{frametitle}{bg=qabg, fg=qatext}
\setbeamercolor{progress bar}{fg=qaaccent, bg=qaline}
\setbeamercolor{alerted text}{fg=qaaccent2}
\setbeamercolor{block title}{bg=qapanel, fg=qaaccent}
\setbeamercolor{block body}{bg=qapanel, fg=qatext}
\setbeamercolor{block title alerted}{bg=qapanel, fg=qaaccent2}
\setbeamercolor{block body alerted}{bg=qapanel, fg=qatext}
\setbeamercolor{block title example}{bg=qapanelalt, fg=qaaccent}
\setbeamercolor{block body example}{bg=qapanelalt, fg=qatext}
\setbeamercolor{example text}{fg=qaaccent}
\setbeamercolor{itemize item}{fg=qaaccent}
\setbeamercolor{itemize subitem}{fg=qaaccent2}
\setbeamercolor{enumerate item}{fg=qaaccent}
\setbeamercolor{enumerate subitem}{fg=qaaccent2}
\setbeamercolor{caption}{fg=qamuted}
\setbeamercolor{page number in head/foot}{fg=qamuted, bg=qabg}
\setbeamercolor{title}{fg=qatext}
\setbeamercolor{subtitle}{fg=qamuted}
\setbeamercolor{author}{fg=qatext}
\setbeamercolor{institute}{fg=qamuted}
\setbeamercolor{date}{fg=qamuted}
\setbeamercolor{section in toc}{fg=qaaccent}
\setbeamercolor{subsection in toc}{fg=qamuted}
% Sommaire de debut de section : section en cours en violet et gras, les autres en gris
\setbeamercolor{section in toc shaded}{fg=qamuted}
\setbeamercolor{subsection in toc shaded}{fg=qamuted}

\setbeamertemplate{blocks}[rounded][shadow=false]
\setbeamertemplate{itemize item}{\raisebox{0.15em}{\scriptsize\textcolor{qaaccent}{\ensuremath{\blacktriangleright}}}}
\setbeamertemplate{itemize subitem}{\raisebox{0.1em}{\tiny\textcolor{qaaccent2}{\ensuremath{\bullet}}}}

\setbeamerfont{frametitle}{series=\bfseries, size=\Large}
\setbeamerfont{title}{series=\bfseries, size=\Large}
\setbeamerfont{subtitle}{size=\normalsize}
\setbeamerfont{institute}{size=\scriptsize}
\setbeamerfont{date}{size=\scriptsize}
\setbeamerfont{author}{size=\small}

\metroset{sectionpage=none, subsectionpage=none, numbering=fraction, progressbar=frametitle}
\setbeamertemplate{section in toc}[sections numbered]

% ---------------------------------------------------------------
% LISTINGS -- theme sombre (lstlisting reste utilisable tel quel)
% \lstcode : police + couleur de base, a reutiliser dans les
%            surcharges locales (ex. basicstyle=\lstcode\scriptsize)
% ---------------------------------------------------------------
\newcommand{\lstcode}{\ttfamily\color{qatext}}

\lstdefinelanguage{TypeScript}{
  keywords={typeof, new, true, false, catch, function, return, null, switch, var,
    if, in, while, do, else, case, break, const, let, interface, class, extends,
    implements, enum, export, import, this, super, constructor, public, private,
    protected, static, async, await, throw, throws, try, finally, void},
  ndkeywords={boolean, any, number, string, int, long, double},
  sensitive=true,
  comment=[l]{//},
  morecomment=[s]{/*}{*/},
  morestring=[b]',
  morestring=[b]"
}

\lstset{
  language=Java,
  basicstyle=\lstcode\footnotesize,
  keywordstyle=\color{qacodekw}\bfseries,
  ndkeywordstyle=\color{qaaccent2}\bfseries,
  identifierstyle=\color{qatext},
  stringstyle=\color{qacodestr},
  commentstyle=\color{qacodecom}\itshape,
  morecomment=[s][\color{qacodeann}]{/**}{*/},
  backgroundcolor=\color{qapanelalt},
  rulecolor=\color{qaline},
  frame=single,
  framesep=3pt,
  xleftmargin=4pt,
  xrightmargin=4pt,
  keepspaces=true,
  columns=fullflexible,
  showstringspaces=false,
  breaklines=true,
  tabsize=2,
  escapeinside={(*@}{@*)},
  % accents UTF-8 dans le code (listings ne gere pas l'UTF-8 multi-octets)
  literate={é}{{\'e}}1 {è}{{\`e}}1 {ê}{{\^e}}1 {à}{{\`a}}1 {â}{{\^a}}1
           {ù}{{\`u}}1 {û}{{\^u}}1 {ô}{{\^o}}1 {î}{{\^i}}1 {ç}{{\c c}}1
           {É}{{\'E}}1 {À}{{\`A}}1 {€}{{\euro{}}}1 {…}{{\ldots}}1,
}

% ---------------------------------------------------------------
% CARD-STYLE BOXES
% ---------------------------------------------------------------
\tcbset{
  figcard/.style={colback=white, colframe=white, boxrule=0pt, arc=4pt,
    left=3pt, right=3pt, top=3pt, bottom=3pt, boxsep=0pt,
    halign=center, valign=center},
  datacard/.style={colback=qapanel, colframe=qaline, boxrule=0.4pt, arc=3pt,
    left=8pt, right=8pt, top=6pt, bottom=6pt, coltext=qatext,
    before upper={\color{qatext}}, halign=center},
  codebox/.style={colback=qapanelalt, colframe=qaline, boxrule=0.4pt, arc=3pt,
    left=8pt, right=8pt, top=5pt, bottom=5pt, listing only,
    listing options={basicstyle=\ttfamily\small\color{qaaccent},
      breaklines=true, showstringspaces=false, columns=fullflexible,
      keepspaces=true, frame=none, backgroundcolor={}}}
}
\newtcolorbox{figcard}[1][]{figcard,#1}
\newtcolorbox{datacard}{datacard}
\newtcblisting{codebox}{codebox}

% ---------------------------------------------------------------
% HANDY MACROS
% ---------------------------------------------------------------
% Mention d'elaboration (discrete : plus sombre que le texte secondaire)
\definecolor{qaaimention}{HTML}{4E5566}
\newcommand{\raimention}{{\tiny\color{qaaimention}Supports \'elabor\'es avec l'aide de Claude Code Opus 5.5}}
\newcommand{\hl}[1]{\alert{#1}}
\newcommand{\source}[1]{{\scriptsize\color{qamuted}#1}}

% Sommaire automatique en debut de section
\AtBeginSection[]{%
  \begin{frame}<beamer>{Plan du cours}
    \tableofcontents[currentsection]
  \end{frame}}

% Page de titre : \rtitleslide{numero du cours}{titre du cours}
% Les auteurs sont repris de \rauthors (modifiable avant \begin{document})
\newcommand{\rauthors}{Chouki Tibermacine}
\newcommand{\rtitleslide}[2]{%
\begin{frame}[plain,noframenumbering]
  \begin{center}
    \vspace{2mm}
    {\color{qaaccent}\rule{2.1cm}{1.2pt}}\\[6mm]
    {\Large\bfseries\color{qatext} R4.02}\\[1mm]
    {\Large\bfseries\color{qatext} Qualit\'e de d\'eveloppement}\\[5mm]
    {\large\color{qamuted} Cours #1}\\[1mm]
    {\large\color{qaaccent} #2}\\[6mm]
    \includegraphics[height=2.1cm]{iut_logo_transparent.png}

    \vspace{4mm}

    {\small\color{qatext} \rauthors}\\[1mm]
    {\scriptsize\color{qamuted} BUT 2 Informatique --- IUT de Vannes --- Universit\'e Bretagne Sud}\\[2mm]
    \raimention
  \end{center}
\end{frame}
}

% ---- Diapo de demonstration : \rdemo{numero}{titre}{contenu (items)}{fichier(s) a ouvrir} ----
\newcommand{\rdemo}[4]{%
\begin{frame}{D\'emo #1 --- #2}
  \begin{tikzpicture}[remember picture, overlay]
    \node[anchor=north east, fill=qaaccent2, text=qabg, font=\scriptsize\bfseries,
      rounded corners=3pt, inner sep=4pt] at ([xshift=-0.4cm, yshift=-0.35cm]current page.north east) {D\'EMO};
  \end{tikzpicture}
  \begin{columns}[T]
    \begin{column}{0.62\textwidth}
      \begin{itemize}
        #3
      \end{itemize}
    \end{column}
    \begin{column}{0.34\textwidth}
      \begin{datacard}
        \raggedright\scriptsize
        {\color{qamuted} Dans l'IDE}\\[1mm]
        #4
      \end{datacard}
    \end{column}
  \end{columns}
\end{frame}}

% ---- Diapo de fin « Questions ? » ----
\newcommand{\rcontact}{chouki.tibermacine@univ-ubs.fr}
\newcommand{\rquestionslide}{%
\begin{frame}{Questions ?}
  \begin{center}
    \vspace{2mm}
    \begin{tikzpicture}
      \node[regular polygon, regular polygon sides=6, draw=qaline, thick,
        minimum size=4.6cm, rotate=90] (hexin) {};
      \node[regular polygon, regular polygon sides=6, draw=qaaccent, dashed,
        thick, minimum size=6.2cm, rotate=90] (hexout) {};
      \node[font=\fontsize{92}{92}\selectfont\bfseries, color=qaaccent2] at (hexin) {?};
    \end{tikzpicture}\\[5mm]
    {\small\color{qamuted} \rcontact}
  \end{center}
\end{frame}}

%   \graphicspath{{../../latex/common/assets/}{img/}{figs/}}
% ---------------------------------------------------------------
\usetheme{metropolis}
\usepackage{appendixnumberbeamer}

\usepackage[T1]{fontenc}
\usepackage[utf8]{inputenc}
\usepackage[french]{babel}
\usepackage[scale=0.94]{tgheros}
\renewcommand{\familydefault}{\sfdefault}
\usepackage{amsmath,amssymb,amsfonts}
\usepackage{graphicx}
\usepackage{booktabs}
\usepackage{multirow}
\usepackage{tikz}
\usetikzlibrary{shapes.geometric, shapes, arrows, arrows.meta, positioning, calc, fit}
\usepackage{xcolor}
\usepackage{listings}
\usepackage[skins,listings]{tcolorbox}
\usepackage{qrcode}
\usepackage{eurosym}
\usepackage{hyperref}

% ---------------------------------------------------------------
% PALETTE (dark-mode tech: near-black canvas, violet (logo IUT) + gold accents)
% ---------------------------------------------------------------
\definecolor{qabg}{HTML}{0E1016}      % canvas background
\definecolor{qapanel}{HTML}{181B24}   % card / block panel
\definecolor{qapanelalt}{HTML}{20242F}% secondary panel tint
\definecolor{qaline}{HTML}{2B303D}    % hairline rules / table rules
\definecolor{qatext}{HTML}{ECEEF2}    % primary text (off-white)
\definecolor{qamuted}{HTML}{8B93A6}   % secondary / caption text
\definecolor{qaaccent}{HTML}{D93ED7}  % primary accent -- violet du logo IUT, eclairci pour fond sombre
\definecolor{qaaccent2}{HTML}{F2B84B} % secondary accent -- or/ambre
\definecolor{qacodekw}{HTML}{C792EA}  % code : mots-cles
\definecolor{qacodestr}{HTML}{C3E88D} % code : chaines
\definecolor{qacodecom}{HTML}{7F8AA3} % code : commentaires
\definecolor{qacodeann}{HTML}{F2B84B} % code : annotations / javadoc

\setbeamercolor{background canvas}{bg=qabg}
\setbeamercolor{normal text}{fg=qatext}
\setbeamercolor{palette primary}{bg=qabg, fg=qamuted}
\setbeamercolor{frametitle}{bg=qabg, fg=qatext}
\setbeamercolor{progress bar}{fg=qaaccent, bg=qaline}
\setbeamercolor{alerted text}{fg=qaaccent2}
\setbeamercolor{block title}{bg=qapanel, fg=qaaccent}
\setbeamercolor{block body}{bg=qapanel, fg=qatext}
\setbeamercolor{block title alerted}{bg=qapanel, fg=qaaccent2}
\setbeamercolor{block body alerted}{bg=qapanel, fg=qatext}
\setbeamercolor{block title example}{bg=qapanelalt, fg=qaaccent}
\setbeamercolor{block body example}{bg=qapanelalt, fg=qatext}
\setbeamercolor{example text}{fg=qaaccent}
\setbeamercolor{itemize item}{fg=qaaccent}
\setbeamercolor{itemize subitem}{fg=qaaccent2}
\setbeamercolor{enumerate item}{fg=qaaccent}
\setbeamercolor{enumerate subitem}{fg=qaaccent2}
\setbeamercolor{caption}{fg=qamuted}
\setbeamercolor{page number in head/foot}{fg=qamuted, bg=qabg}
\setbeamercolor{title}{fg=qatext}
\setbeamercolor{subtitle}{fg=qamuted}
\setbeamercolor{author}{fg=qatext}
\setbeamercolor{institute}{fg=qamuted}
\setbeamercolor{date}{fg=qamuted}
\setbeamercolor{section in toc}{fg=qaaccent}
\setbeamercolor{subsection in toc}{fg=qamuted}
% Sommaire de debut de section : section en cours en violet et gras, les autres en gris
\setbeamercolor{section in toc shaded}{fg=qamuted}
\setbeamercolor{subsection in toc shaded}{fg=qamuted}

\setbeamertemplate{blocks}[rounded][shadow=false]
\setbeamertemplate{itemize item}{\raisebox{0.15em}{\scriptsize\textcolor{qaaccent}{\ensuremath{\blacktriangleright}}}}
\setbeamertemplate{itemize subitem}{\raisebox{0.1em}{\tiny\textcolor{qaaccent2}{\ensuremath{\bullet}}}}

\setbeamerfont{frametitle}{series=\bfseries, size=\Large}
\setbeamerfont{title}{series=\bfseries, size=\Large}
\setbeamerfont{subtitle}{size=\normalsize}
\setbeamerfont{institute}{size=\scriptsize}
\setbeamerfont{date}{size=\scriptsize}
\setbeamerfont{author}{size=\small}

\metroset{sectionpage=none, subsectionpage=none, numbering=fraction, progressbar=frametitle}
\setbeamertemplate{section in toc}[sections numbered]

% ---------------------------------------------------------------
% LISTINGS -- theme sombre (lstlisting reste utilisable tel quel)
% \lstcode : police + couleur de base, a reutiliser dans les
%            surcharges locales (ex. basicstyle=\lstcode\scriptsize)
% ---------------------------------------------------------------
\newcommand{\lstcode}{\ttfamily\color{qatext}}

\lstdefinelanguage{TypeScript}{
  keywords={typeof, new, true, false, catch, function, return, null, switch, var,
    if, in, while, do, else, case, break, const, let, interface, class, extends,
    implements, enum, export, import, this, super, constructor, public, private,
    protected, static, async, await, throw, throws, try, finally, void},
  ndkeywords={boolean, any, number, string, int, long, double},
  sensitive=true,
  comment=[l]{//},
  morecomment=[s]{/*}{*/},
  morestring=[b]',
  morestring=[b]"
}

\lstset{
  language=Java,
  basicstyle=\lstcode\footnotesize,
  keywordstyle=\color{qacodekw}\bfseries,
  ndkeywordstyle=\color{qaaccent2}\bfseries,
  identifierstyle=\color{qatext},
  stringstyle=\color{qacodestr},
  commentstyle=\color{qacodecom}\itshape,
  morecomment=[s][\color{qacodeann}]{/**}{*/},
  backgroundcolor=\color{qapanelalt},
  rulecolor=\color{qaline},
  frame=single,
  framesep=3pt,
  xleftmargin=4pt,
  xrightmargin=4pt,
  keepspaces=true,
  columns=fullflexible,
  showstringspaces=false,
  breaklines=true,
  tabsize=2,
  escapeinside={(*@}{@*)},
  % accents UTF-8 dans le code (listings ne gere pas l'UTF-8 multi-octets)
  literate={é}{{\'e}}1 {è}{{\`e}}1 {ê}{{\^e}}1 {à}{{\`a}}1 {â}{{\^a}}1
           {ù}{{\`u}}1 {û}{{\^u}}1 {ô}{{\^o}}1 {î}{{\^i}}1 {ç}{{\c c}}1
           {É}{{\'E}}1 {À}{{\`A}}1 {€}{{\euro{}}}1 {…}{{\ldots}}1,
}

% ---------------------------------------------------------------
% CARD-STYLE BOXES
% ---------------------------------------------------------------
\tcbset{
  figcard/.style={colback=white, colframe=white, boxrule=0pt, arc=4pt,
    left=3pt, right=3pt, top=3pt, bottom=3pt, boxsep=0pt,
    halign=center, valign=center},
  datacard/.style={colback=qapanel, colframe=qaline, boxrule=0.4pt, arc=3pt,
    left=8pt, right=8pt, top=6pt, bottom=6pt, coltext=qatext,
    before upper={\color{qatext}}, halign=center},
  codebox/.style={colback=qapanelalt, colframe=qaline, boxrule=0.4pt, arc=3pt,
    left=8pt, right=8pt, top=5pt, bottom=5pt, listing only,
    listing options={basicstyle=\ttfamily\small\color{qaaccent},
      breaklines=true, showstringspaces=false, columns=fullflexible,
      keepspaces=true, frame=none, backgroundcolor={}}}
}
\newtcolorbox{figcard}[1][]{figcard,#1}
\newtcolorbox{datacard}{datacard}
\newtcblisting{codebox}{codebox}

% ---------------------------------------------------------------
% HANDY MACROS
% ---------------------------------------------------------------
% Mention d'elaboration (discrete : plus sombre que le texte secondaire)
\definecolor{qaaimention}{HTML}{4E5566}
\newcommand{\raimention}{{\tiny\color{qaaimention}Supports \'elabor\'es avec l'aide de Claude Code Opus 5.5}}
\newcommand{\hl}[1]{\alert{#1}}
\newcommand{\source}[1]{{\scriptsize\color{qamuted}#1}}

% Sommaire automatique en debut de section
\AtBeginSection[]{%
  \begin{frame}<beamer>{Plan du cours}
    \tableofcontents[currentsection]
  \end{frame}}

% Page de titre : \rtitleslide{numero du cours}{titre du cours}
% Les auteurs sont repris de \rauthors (modifiable avant \begin{document})
\newcommand{\rauthors}{Chouki Tibermacine}
\newcommand{\rtitleslide}[2]{%
\begin{frame}[plain,noframenumbering]
  \begin{center}
    \vspace{2mm}
    {\color{qaaccent}\rule{2.1cm}{1.2pt}}\\[6mm]
    {\Large\bfseries\color{qatext} R4.02}\\[1mm]
    {\Large\bfseries\color{qatext} Qualit\'e de d\'eveloppement}\\[5mm]
    {\large\color{qamuted} Cours #1}\\[1mm]
    {\large\color{qaaccent} #2}\\[6mm]
    \includegraphics[height=2.1cm]{iut_logo_transparent.png}

    \vspace{4mm}

    {\small\color{qatext} \rauthors}\\[1mm]
    {\scriptsize\color{qamuted} BUT 2 Informatique --- IUT de Vannes --- Universit\'e Bretagne Sud}\\[2mm]
    \raimention
  \end{center}
\end{frame}
}

% ---- Diapo de demonstration : \rdemo{numero}{titre}{contenu (items)}{fichier(s) a ouvrir} ----
\newcommand{\rdemo}[4]{%
\begin{frame}{D\'emo #1 --- #2}
  \begin{tikzpicture}[remember picture, overlay]
    \node[anchor=north east, fill=qaaccent2, text=qabg, font=\scriptsize\bfseries,
      rounded corners=3pt, inner sep=4pt] at ([xshift=-0.4cm, yshift=-0.35cm]current page.north east) {D\'EMO};
  \end{tikzpicture}
  \begin{columns}[T]
    \begin{column}{0.62\textwidth}
      \begin{itemize}
        #3
      \end{itemize}
    \end{column}
    \begin{column}{0.34\textwidth}
      \begin{datacard}
        \raggedright\scriptsize
        {\color{qamuted} Dans l'IDE}\\[1mm]
        #4
      \end{datacard}
    \end{column}
  \end{columns}
\end{frame}}

% ---- Diapo de fin « Questions ? » ----
\newcommand{\rcontact}{chouki.tibermacine@univ-ubs.fr}
\newcommand{\rquestionslide}{%
\begin{frame}{Questions ?}
  \begin{center}
    \vspace{2mm}
    \begin{tikzpicture}
      \node[regular polygon, regular polygon sides=6, draw=qaline, thick,
        minimum size=4.6cm, rotate=90] (hexin) {};
      \node[regular polygon, regular polygon sides=6, draw=qaaccent, dashed,
        thick, minimum size=6.2cm, rotate=90] (hexout) {};
      \node[font=\fontsize{92}{92}\selectfont\bfseries, color=qaaccent2] at (hexin) {?};
    \end{tikzpicture}\\[5mm]
    {\small\color{qamuted} \rcontact}
  \end{center}
\end{frame}}

\graphicspath{{../../latex/common/assets/}{img/}}

\title{Environnements, intégration continue et couverture de tests}
\author[C. Tibermacine]{Chouki Tibermacine}
\date{}

\lstdefinelanguage{yaml}{
  keywords={true,false,null},
  sensitive=false,
  comment=[l]{\#},
  morestring=[b]',
  morestring=[b]"
}
\lstdefinestyle{yml}{language=yaml, basicstyle=\lstcode\scriptsize}
\lstdefinestyle{sortie}{language={}, basicstyle=\lstcode\scriptsize, keywordstyle={}, commentstyle={}, stringstyle={}}

\tikzset{
  env/.style={rectangle, rounded corners=3pt, draw=qaaccent, thick, fill=qapanel, text=qatext,
    minimum height=0.9cm, minimum width=2.5cm, align=center, font=\scriptsize},
  job/.style={rectangle, rounded corners=2pt, draw=qamuted, thick, fill=qapanelalt, text=qatext,
    minimum height=0.5cm, minimum width=1.9cm, align=center, font=\tiny},
  arr/.style={-Latex, qamuted, thick}
}

\begin{document}

\rtitleslide{6}{Environnements, intégration continue et couverture de tests}

\begin{frame}{Où en est-on ?}
  \begin{columns}[c]
    \begin{column}{0.5\textwidth}
      \begin{itemize}
        \item Nous savons écrire tous les niveaux de tests : unitaires, acceptation, intégration, bout en bout
        \item Trois questions restent :
          \begin{enumerate}
            \item \textbf{Où} s'exécute l'application, et quels tests dans quel environnement ?
            \item \textbf{Qui} lance les tests, et \textbf{quand} ? (« j'ai oublié de les lancer »)
            \item Nos tests sont-ils \textbf{suffisants} ?
          \end{enumerate}
      \end{itemize}
    \end{column}
    \begin{column}{0.46\textwidth}
      \begin{datacard}
        \small
        Réponses d'aujourd'hui :\\[1mm]
        1. les \textbf{environnements}\\
        2. l'\textbf{intégration continue} (GitLab CI)\\
        3. la \textbf{couverture} (JaCoCo) et la \textbf{mutation} (PIT)
      \end{datacard}
    \end{column}
  \end{columns}
\end{frame}

% ================================================================
\section{Les environnements d'une application}

\begin{frame}{Du poste du développeur à la production}
  \begin{center}
  \begin{tikzpicture}[node distance=0.55cm]
    \node[env] (dev) {Développement\\{\tiny\color{qamuted}votre poste}};
    \node[env, right=of dev] (ci) {Intégration\\{\tiny\color{qamuted}machines de CI}};
    \node[env, right=of ci] (rec) {Recette\\{\tiny\color{qamuted}\textit{staging}, préproduction}};
    \node[env, right=of rec, draw=qaaccent2] (prod) {Production\\{\tiny\color{qamuted}les vrais clients}};
    \draw[arr] (dev) -- node[above, font=\tiny, text=qamuted] {\textit{push}} (ci);
    \draw[arr] (ci) -- node[above, font=\tiny, text=qamuted] {auto.} (rec);
    \draw[arr] (rec) -- node[above, font=\tiny, text=qamuted] {décision} (prod);
  \end{tikzpicture}
  \end{center}
  \vspace{2mm}
  \begin{itemize}
    \item \textbf{Développement} : on code, on lance les tests rapides, base locale (\texttt{docker compose})
    \item \textbf{Intégration} : à chaque \textit{push}, une machine neutre compile et lance \textbf{tous} les tests
    \item \textbf{Recette} : une copie aussi fidèle que possible de la production ; le client y valide les stories (CM2)
    \item \textbf{Production} : ce que les clients utilisent ; on n'y fait pas d'essais
  \end{itemize}
\end{frame}

\begin{frame}{Quels tests dans quel environnement ?}
  \small
  \begin{datacard}
    \begin{tabular}{@{}p{2.4cm}p{10.2cm}@{}}
      \toprule
      \textbf{Environnement} & \textbf{Tests} \\
      \midrule
      Développement & unitaires (secondes, à chaque modification), quelques tests d'intégration ciblés \\[0.5mm]
      Intégration (CI) & \textbf{tous} les tests automatiques : unitaires, intégration (Testcontainers), acceptation (Cucumber), bout en bout (Playwright) ; couverture, mutation \\[0.5mm]
      Recette & bout en bout sur l'application déployée ; \textbf{recette} par le client ; tests de performance ; tests d'utilisabilité \\[0.5mm]
      Production & \textit{smoke tests} juste après le déploiement (« l'application répond, la connexion marche ») ; supervision : erreurs, temps de réponse \\
      \bottomrule
    \end{tabular}
  \end{datacard}
  \vspace{2mm}
  \begin{itemize}
    \item Plus on avance vers la production, plus un bug coûte cher (CM1) : chaque environnement est un \textbf{filet} de plus
  \end{itemize}
\end{frame}

\begin{frame}{Construire une fois, configurer partout}
  \begin{columns}[T]
    \begin{column}{0.48\textwidth}
      \begin{block}{Le même artefact}
        \begin{itemize}
          \item On construit \textbf{une seule fois} (un \texttt{.jar}, une image Docker), puis on \textbf{promeut} ce même artefact de la CI à la recette, puis à la production
          \item Ce qui a été testé est exactement ce qui est déployé
          \item Seule la \textbf{configuration} change : URL de la base, adresses des services, niveau de journalisation (profils Spring, variables d'environnement)
        \end{itemize}
      \end{block}
    \end{column}
    \begin{column}{0.48\textwidth}
      \begin{alertblock}{Secrets et données}
        \begin{itemize}
          \item Mots de passe et clés d'API : \textbf{jamais} dans le dépôt ; variables protégées de la CI, coffre-fort de secrets
          \item Développement et CI : jeux de données de test
          \item Recette : données réalistes mais \textbf{anonymisées} (RGPD)
          \item Production : les vraies données ; jamais de tests qui les modifient
        \end{itemize}
      \end{alertblock}
    \end{column}
  \end{columns}
\end{frame}

% ================================================================
\section{Intégration et déploiement continus}

\begin{frame}{CI, livraison continue, déploiement continu}
  \begin{columns}[T]
    \begin{column}{0.5\textwidth}
      \begin{itemize}
        \item \textbf{Intégration continue} (CI) : chaque modification poussée est compilée et testée automatiquement, plusieurs fois par jour
        \item \textbf{Livraison continue} : chaque version qui passe la CI est \textbf{prête} à être déployée ; la mise en production est une décision humaine (un bouton)
        \item \textbf{Déploiement continu} : la mise en production est elle aussi automatique
      \end{itemize}
    \end{column}
    \begin{column}{0.46\textwidth}
      \begin{block}{La non-régression, automatisée}
        \begin{itemize}
          \item Toute la suite est rejouée à chaque \textit{push} et à chaque \textit{merge request}
          \item Une \textit{merge request} dont la pipeline est rouge \textbf{ne peut pas} être fusionnée
          \item La règle d'or du CM1 devient vérifiable : un bug corrigé $=$ un test ajouté, rejoué à jamais
          \item Personne ne peut « oublier » de lancer les tests
        \end{itemize}
      \end{block}
    \end{column}
  \end{columns}
\end{frame}

\begin{frame}{Anatomie d'une pipeline GitLab CI}
  \begin{center}
  \begin{tikzpicture}[x=2.15cm, y=0.75cm]
    \foreach \x/\s in {0/build, 1/test, 2/e2e, 3/qualite, 4/recette, 5/production} {
      \node[font=\scriptsize\bfseries, text=qaaccent] at (\x,1.1) {\s};
      \draw[qaline] (\x-0.47,0.75) rectangle (\x+0.47,-1.65);
    }
    \node[job] at (0,0) {compiler};
    \node[job] at (1,0) {tests};
    \node[job] at (2,0) {e2e};
    \node[job] at (3,0) {mutation};
    \node[job] at (4,0) {recette};
    \node[job, draw=qaaccent2] at (5,0) {production\\{\tiny manuel}};
    \foreach \x in {0,...,4} \draw[arr] (\x+0.47,0) -- (\x+0.53,0);
    \node[font=\tiny, text=qamuted, align=center] at (1,-1.0) {service\\\texttt{docker:dind}};
    \node[font=\tiny, text=qamuted, align=center] at (2,-1.0) {image avec\\navigateurs};
    \node[font=\tiny, text=qamuted, align=center] at (3,-1.0) {branche\\principale};
  \end{tikzpicture}
  \end{center}
  \small
  \begin{itemize}
    \item Fichier \texttt{.gitlab-ci.yml} à la racine du dépôt, versionné avec le code
    \item \textbf{Étapes} (\textit{stages}) exécutées dans l'ordre ; les \textbf{jobs} d'une même étape en parallèle ; si un job échoue, les étapes suivantes ne sont pas lancées
    \item Chaque job s'exécute dans un conteneur neuf (\texttt{image}), sur une machine appelée \textit{runner}
    \item \textbf{Artefacts} : fichiers conservés après le job (rapports, traces) ; \textbf{cache} : dépendances Maven réutilisées d'une exécution à l'autre
  \end{itemize}
\end{frame}

\begin{frame}[fragile]{Le fichier \texttt{.gitlab-ci.yml} (1/2)}
  \begin{columns}[T]
    \begin{column}{0.6\textwidth}
\begin{lstlisting}[style=yml]
stages: [build, test, e2e, qualite, recette, production]
default:
  image: maven:3.9-eclipse-temurin-21
  cache: { paths: [.m2/repository] }
variables:
  DOCKER_HOST: "tcp://docker:2375"   # pour Testcontainers
  DOCKER_TLS_CERTDIR: ""
tests:
  stage: test
  services:
    - { name: "docker:dind", command: ["--tls=false"] }
  script:
    - mvn --batch-mode verify -Dgroups='!e2e'
  artifacts:
    when: always
    reports:
      junit: [target/surefire-reports/TEST-*.xml]
      coverage_report:
        coverage_format: jacoco
        path: target/site/jacoco/jacoco.xml
\end{lstlisting}
    \end{column}
    \begin{column}{0.37\textwidth}
      \small
      \begin{itemize}
        \item \texttt{image} : Java et Maven, sans rien installer
        \item \texttt{services} : un Docker dans le job, pour que Testcontainers démarre PostgreSQL (le \textit{runner} doit l'autoriser)
        \item \texttt{-Dgroups='!e2e'} : tous les tests sauf ceux tagués \texttt{e2e}, lancés plus tard
        \item \texttt{reports} : GitLab lit les résultats JUnit et la couverture et les affiche dans la \textit{merge request}
        \item \texttt{when: always} : garder les rapports même si un test échoue
      \end{itemize}
    \end{column}
  \end{columns}
\end{frame}

\begin{frame}[fragile]{Le fichier \texttt{.gitlab-ci.yml} (2/2)}
  \begin{columns}[T]
    \begin{column}{0.6\textwidth}
\begin{lstlisting}[style=yml]
e2e:
  stage: e2e
  image: mcr.microsoft.com/playwright/java:v1.64.0-noble
  services:
    - { name: "docker:dind", command: ["--tls=false"] }
  script:
    - mvn --batch-mode verify -Dgroups=e2e
  artifacts: { when: on_failure, paths: [target/traces/] }
recette:
  stage: recette
  script: echo "Déploiement de $CI_COMMIT_SHORT_SHA"
  environment: { name: recette }
production:
  stage: production
  script: echo "Déploiement de $CI_COMMIT_SHORT_SHA"
  environment: { name: production }
  when: manual
\end{lstlisting}
    \end{column}
    \begin{column}{0.37\textwidth}
      \small
      \begin{itemize}
        \item Image officielle Playwright : Java, Maven et navigateurs déjà installés
        \item En cas d'échec, les \textbf{traces} Playwright sont téléchargeables depuis la pipeline (CM5)
        \item \texttt{environment} : GitLab garde l'historique de ce qui a été déployé où
        \item \texttt{when: manual} : la mise en production attend un clic (livraison continue)
        \item Ici, le déploiement est simulé (\texttt{echo}) : il faudrait un serveur cible
      \end{itemize}
    \end{column}
  \end{columns}
\end{frame}

\begin{frame}{Organiser la suite de tests}
  \begin{columns}[T]
    \begin{column}{0.48\textwidth}
      \begin{block}{Échouer vite}
        \begin{itemize}
          \item Les tests les plus \textbf{rapides} d'abord : une erreur de compilation ou un test unitaire rouge arrête tout en une minute
          \item Les tests lents (bout en bout, mutation) ensuite, ou seulement sur la branche principale
          \item Une pipeline qui dure 40 minutes, personne n'attend son résultat
        \end{itemize}
      \end{block}
    \end{column}
    \begin{column}{0.48\textwidth}
      \begin{block}{Séparer les tests}
        \begin{itemize}
          \item \textbf{Tags} JUnit (\texttt{@Tag("e2e")}) et filtres Maven (\texttt{-Dgroups=e2e}, \texttt{-Dgroups='!e2e'})
          \item Ou conventions de nommage : \texttt{*Test} pour les tests unitaires (\textit{Surefire}, \texttt{mvn test}), \texttt{*IT} pour l'intégration (\textit{Failsafe}, \texttt{mvn verify})
          \item Un test instable bloque toute l'équipe : on le répare tout de suite (CM5)
        \end{itemize}
      \end{block}
    \end{column}
  \end{columns}
\end{frame}

% ================================================================
\section{Couverture et mutation}

\begin{frame}[fragile]{Mesurer la couverture}
  \begin{columns}[T]
    \begin{column}{0.5\textwidth}
\begin{lstlisting}[basicstyle=\lstcode\scriptsize]
public BigDecimal fraisVirement(BigDecimal montant,
        boolean instantane, boolean premium) {
    if (premium) {
        return BigDecimal.ZERO;
    }
    if (!instantane) {
        return BigDecimal.ZERO;
    }
    if (montant.compareTo(SEUIL) > 0) {
        return FRAIS_INSTANTANE_ELEVE;   // 2,00 €
    }
    return FRAIS_INSTANTANE;             // 1,00 €
}
\end{lstlisting}
    \end{column}
    \begin{column}{0.46\textwidth}
      \small
      \begin{itemize}
        \item \textbf{Couverture des lignes} : proportion des lignes exécutées par au moins un test
        \item \textbf{Couverture des branches} : chaque \texttt{if} a deux branches (vrai, faux) ; les a-t-on toutes parcourues ?
        \item Ici, deux tests (premium, standard) exécutent 5 lignes sur 9\ldots{} et 3 branches sur 6 : les virements instantanés ne sont \textbf{jamais} testés
        \item La couverture des branches est plus exigeante, donc plus utile
      \end{itemize}
    \end{column}
  \end{columns}
\end{frame}

\begin{frame}[fragile]{JaCoCo : rapport et seuil}
  \begin{columns}[T]
    \begin{column}{0.5\textwidth}
\begin{lstlisting}[language=XML, basicstyle=\lstcode\scriptsize]
<plugin>
  <groupId>org.jacoco</groupId>
  <artifactId>jacoco-maven-plugin</artifactId>
  <!-- executions : prepare-agent,
       report, check (phase verify) -->
  <configuration>
    <rules><rule>
      <element>BUNDLE</element>
      <limits><limit>
        <counter>BRANCH</counter>
        <value>COVEREDRATIO</value>
        <minimum>0.80</minimum>
      </limit></limits>
    </rule></rules>
  </configuration>
</plugin>
\end{lstlisting}
    \end{column}
    \begin{column}{0.46\textwidth}
      \small
      \begin{itemize}
        \item JaCoCo observe l'exécution des tests et produit un rapport HTML : lignes \textcolor{qacodestr}{vertes} (couvertes), \textcolor{qaaccent2}{jaunes} (branches en partie couvertes), \textcolor{qaaccent}{rouges} (jamais exécutées)
        \item \texttt{check} : le build \textbf{échoue} sous le seuil
      \end{itemize}
\begin{lstlisting}[style=sortie]
Rule violated for bundle: branches covered
ratio is 0.50, but expected minimum is 0.80
BUILD FAILURE
\end{lstlisting}
      \begin{itemize}
        \item En CI, la couverture s'affiche dans la \textit{merge request}, ligne par ligne
      \end{itemize}
    \end{column}
  \end{columns}
\end{frame}

\begin{frame}[fragile]{Ce que la couverture ne dit pas}
  \begin{columns}[T]
    \begin{column}{0.5\textwidth}
\begin{lstlisting}[basicstyle=\lstcode\scriptsize]
@Test
void instantane_petitMontant() {
    assertThat(tarification
        .fraisVirement(eur("50"), true, false))
        .isPositive();
}
@Test
void instantane_grosMontant() {
    assertThat(tarification
        .fraisVirement(eur("5000"), true, false))
        .isPositive();
}
\end{lstlisting}
      \small Avec deux tests de ce genre en plus : \textbf{100 \%} des lignes et des branches.
    \end{column}
    \begin{column}{0.46\textwidth}
      \small
      \begin{itemize}
        \item Le code est \textbf{exécuté}, mais est-il \textbf{vérifié} ?
          \begin{itemize}
            \item \texttt{isPositive()} accepte 1 \euro{} comme 2 \euro{}
            \item personne ne teste exactement 1\,000 \euro{} : \texttt{>} ou \texttt{>=}, aucun test ne fait la différence
          \end{itemize}
        \item Un test sans assertion « couvre » aussi tout ce qu'il appelle
        \item 100 \% de couverture ne garantit pas l'absence de bugs ; 40 \% garantit que la majorité du code n'est jamais testée
        \item Un seuil imposé devient un objectif : on écrit des tests pour le chiffre, pas pour trouver des bugs
      \end{itemize}
    \end{column}
  \end{columns}
\end{frame}

\begin{frame}{Les tests de mutation (PIT)}
  \begin{columns}[T]
    \begin{column}{0.5\textwidth}
      \begin{itemize}
        \item Idée : \textbf{introduire des bugs} exprès, et vérifier que les tests les détectent
        \item PIT crée des \textbf{mutants} : des copies du code avec une petite modification
          \begin{itemize}
            \item \texttt{>} devient \texttt{>=} (\textit{conditional boundary})
            \item une condition est inversée
            \item une valeur renvoyée est remplacée par \texttt{null}
          \end{itemize}
        \item Il relance les tests sur chaque mutant :
          \begin{itemize}
            \item un test échoue : mutant \textbf{tué} (bien)
            \item tous passent : mutant \alert{survivant}, un bug que nos tests laisseraient passer
          \end{itemize}
      \end{itemize}
    \end{column}
    \begin{column}{0.46\textwidth}
      \begin{datacard}
        \small
        \begin{tabular}{@{}lcc@{}}
          \toprule
          \textbf{Tests} & \textbf{Lignes} & \textbf{Mutants tués} \\
          \midrule
          faibles & 100 \% & \alert{75 \%} (6/8) \\
          solides & 100 \% & 100 \% (8/8) \\
          \bottomrule
        \end{tabular}
      \end{datacard}
      \small
      \begin{itemize}
        \item Même couverture, qualité très différente : le \textbf{score de mutation} mesure la force des assertions
        \item Coût : les tests sont relancés des dizaines de fois ; on le lance sur la branche principale, ou sur le code modifié
      \end{itemize}
    \end{column}
  \end{columns}
\end{frame}

\rdemo{1}{Couverture et mutation}{%
  \item Deux tests seulement : JaCoCo fait \textbf{échouer le build} (50 \% des branches) ; le rapport HTML montre en rouge ce qui n'est jamais exécuté
  \item Des tests faibles : 100 \% de couverture, build vert\ldots
  \item \ldots{} mais PIT : 2 mutants survivent, dont \texttt{>} remplacé par \texttt{>=}
  \item Les tests solides (valeurs exactes, valeur limite 1\,000 \euro{}) : tous les mutants sont tués
  \item Le fichier \texttt{.gitlab-ci.yml} complet
}{\texttt{Tarification}\\\texttt{Tarification*Test}\\\texttt{target/site/jacoco/}\\\texttt{target/pit-reports/}\\\texttt{gitlab-ci-exemple.yml}}

\begin{frame}{Bilan du module : quel test, pour quel risque, où ?}
  \footnotesize
  \begin{datacard}
    \begin{tabular}{@{}p{2.5cm}p{4.1cm}p{2.7cm}p{2.9cm}@{}}
      \toprule
      \textbf{Test} & \textbf{Risque couvert} & \textbf{Outils} & \textbf{Où, quand} \\
      \midrule
      Unitaire & règle métier fausse, valeur limite & JUnit, AssertJ, Mockito & poste, CI ; chaque modification \\[0.5mm]
      Acceptation & on construit le mauvais produit & Cucumber & CI ; recette \\[0.5mm]
      Intégration & SQL, mapping, transaction, configuration & Testcontainers, Spring, MockMvc & CI \\[0.5mm]
      Bout en bout & parcours utilisateur cassé, bug du front & Playwright & CI, recette \\[0.5mm]
      Non-régression & ce qui marchait ne marche plus & toute la suite & CI, à chaque \textit{push} \\[0.5mm]
      Couverture, mutation & tests insuffisants ou trop faibles & JaCoCo, PIT & CI, branche principale \\
      \bottomrule
    \end{tabular}
  \end{datacard}
  \vspace{1mm}
  \small Pour aller plus loin : analyse statique du code, \textit{quality gates} (bloquer une fusion selon des critères de qualité), analyse de sécurité des dépendances et des images dans la pipeline.
\end{frame}

\begin{frame}{À retenir}
  \footnotesize
  \begin{itemize}
    \item Quatre environnements : développement, intégration, \textbf{recette}, production ; un filet de tests à chaque étape
    \item On construit \textbf{une fois} et on promeut le même artefact ; seule la configuration change ; jamais de secret dans le dépôt
    \item \textbf{Intégration continue} : tous les tests, à chaque \textit{push} ; une \textit{merge request} rouge ne se fusionne pas
    \item \texttt{.gitlab-ci.yml} : étapes, jobs, image, services (\texttt{docker:dind} pour Testcontainers), artefacts, environnements, \texttt{when: manual}
    \item Échouer vite : tests rapides d'abord, tests lents ensuite
    \item \textbf{Couverture} (JaCoCo) : ce qui est exécuté ; branches plutôt que lignes ; un seuil qui fait échouer le build
    \item \textbf{Mutation} (PIT) : ce qui est vraiment vérifié ; un mutant survivant est un bug que les tests laisseraient passer
  \end{itemize}
  \vspace{1mm}
  \begin{exampleblock}{TP6}
    La pipeline GitLab CI du projet fil rouge, avec couverture et mutation. \alert{Quiz 2} en fin de semaine, lors d'une séance à part (CM4 à CM6).
  \end{exampleblock}
\end{frame}

\rquestionslide

\end{document}
